% Chapter 6: Conclusion
\chapter{Conclusion}

\section{Summary of the Study}

The thesis built and evaluated an end-to-end imputation--clustering framework for severely incomplete, multi-source social data, using newspaper-reported suicide cases from Bangladesh as a worked example. The choice of dataset reflects two facts: Bangladesh has no centralised national suicide registry, and the available newspaper-derived records combine extreme non-random missingness (often above 50--90\% for several important features) with multi-cohort heterogeneity that defeats off-the-shelf machine-learning pipelines.

The original 760-record Kaggle dataset was expanded through manual collection of 857 additional cases from verified Bangladeshi news sources, giving 1,617 records and 34 analytical columns covering 2017--2025. Nine new socioeconomic and clinical attributes were added, and Geopy and Meteostat APIs sharply reduced missingness in spatial and weather fields. Three imputation methods (MICE, KNN, MissForest) were each applied under two predictor-selection regimes, producing six imputed datasets plus two non-imputed baselines. Four clustering algorithms (K-Means, Agglomerative, DBSCAN, Spectral) were then applied across all eight datasets, yielding 32 end-to-end pipelines. Every pipeline was scored against five standard internal validity indices and a study-specific balance penalty included to demote the degenerate-but-tight partitions that single-index validation rewards on imbalanced social data.

The best-performing combination --- Agglomerative Clustering with MissForest no-context imputation --- achieved a composite score of 27.30 out of 32. Its two-cluster solution is the subject of the detailed analysis in Chapter~4 and the discussion below.

\section{Key Findings}

The study produced five findings spanning methodological, empirical, and evaluative dimensions.

Finding 1 --- \textit{MissForest consistently outperformed MICE and KNN.} Across both predictor strategies and all four clustering algorithms, MissForest-imputed datasets produced more consistent, better-ranked solutions than either MICE or KNN. The pattern mirrors the wider imputation-benchmark literature \cite{ref12,ref32,ref33} and reinforces the suitability of tree-based non-parametric imputation for mixed-type data with severe missingness.

Finding 2 --- \textit{Imputed datasets vastly outperformed non-imputed baselines.} Both baseline datasets ranked at or near the bottom of the composite table. For data this incomplete, attempting to cluster without principled imputation produces structurally weak and practically uninterpretable partitions --- a finding with direct implications for future researchers working with similarly incomplete datasets in LMICs.

Finding 3 --- \textit{The balance penalty was load-bearing for fair ranking.} Several DBSCAN configurations posted superficially high Silhouette and Calinski--Harabasz values while flagging over 99\% of points as noise; without the 35\%-weighted balance term (see \S3.6.3 for the formulation) these would have ranked among the leading solutions. The result echoes the broader caution against single-index clustering evaluation on imbalanced data \cite{ref54} and motivates the balance term as a pragmatic, project-level scoring choice rather than a general clustering metric.

Finding 4 --- \textit{The composite framework recovered source-cohort structure as the dominant signal on multi-source newspaper data, a methodologically informative observation.} The two-cluster solution from the best pipeline (agg\_rf\_no\_context) aligns almost perfectly with the dataset's two-cohort construction history: Cluster~0 contains the 759 records inherited from the 2020-focused Kaggle subset, and Cluster~1 contains the 858 manually collected 2017--2019 and 2021--2025 records. The clustering algorithm had no access to collection year, source, or any temporal identifier; the boundary emerged from feature-level patterns alone. Read carefully, the result is both a validation and a caveat: a validation in that the imputation preserved cohort-level differences rather than collapsing them, and a caveat in that the two cohorts differ systematically in completeness, reporting era, and the journalistic and geographic profile of the underlying news coverage, with those differences propagating into virtually every downstream feature. The partition is therefore best characterised as a source-cohort split that co-varies with --- rather than cleanly isolates --- substantive demographic, environmental, and behavioural variation. A secondary implication is that the missingness pattern in this dataset is itself informative: cases with higher missingness rates reflect a distinct reporting environment, consistent with a not-missing-at-random (MNAR) structure that constrains how strongly any clustering result on this dataset can be generalised. The confound is recorded explicitly as Limitation~\ref{lim:source_cohort} in \S6.4, and the cluster profiles in Chapter~4 are framed throughout as exploratory rather than confirmatory.

Finding 5 --- \textit{Spectral Clustering had the strongest algorithm-level average; Agglomerative was the more robust practical choice.} Spectral achieved the highest average performance across the composite metric components, but its $O(n^3)$ cost and sensitivity to parameter choices made it less practically reliable than Agglomerative, which combined second-best average performance with greater stability across imputation variants.

\section{Achievement of Research Objectives}

All six research objectives from Chapter~1 were achieved within the scope of the study. Table~\ref{tbl:research_objectives_and_their} maps each objective to its concrete outcome.

\begin{table}[H]
\centering
\caption[Research objectives and their achievement status]{Research objectives and their achievement status}
\label{tbl:research_objectives_and_their}
\renewcommand{\arraystretch}{1.0}
\small
\begin{tabularx}{\textwidth}{|>{\hsize=1.165\hsize\linewidth=\hsize\raggedright\arraybackslash}X|>{\hsize=1.388\hsize\linewidth=\hsize\raggedright\arraybackslash}X|>{\hsize=0.447\hsize\linewidth=\hsize\raggedright\arraybackslash}X|}
\hline
\textbf{Objective} & \textbf{Achievement} & \textbf{Status} \\ \hline
Expand dataset 760 $\to$ 1,617 records with 9 new attributes (2017--2025) & 1,617 records, 34 analytical columns, 9 added socioeconomic/clinical attributes & Achieved \\ \hline
Reduce missingness in spatial/weather features via Geopy and Meteostat APIs & Lat/long missingness reduced from \textasciitilde{}53\% to 6.49\%; partial reduction in weather variables & Achieved \\ \hline
Compare MICE, KNN, MissForest under context-aware and no-context predictor strategies & Six imputed datasets produced; all post-imputation validation checks passed & Achieved \\ \hline
Apply four clustering algorithms across all 8 datasets (32 pipelines) & All 32 pipelines executed, evaluated, and ranked & Achieved \\ \hline
Evaluate via 5 standard validity indices plus a study-specific balance penalty & Composite scoring framework with balance term at 35\% of total composite score & Achieved \\ \hline
Identify the best pipeline and profile its clusters & agg\_\hspace{0pt}rf\_\hspace{0pt}no\_\hspace{0pt}context selected (27.30/32); two-cluster solution profiled across demographic, geographic, and temporal dimensions & Achieved \\ \hline
\end{tabularx}
\end{table}

Objective~6 deserves a brief qualification: the two-cluster solution is interpretable, but the largely cohort-driven separation means the demographic, environmental, and geographic readings of the clusters are exploratory rather than confirmatory. More homogeneously sourced data would support sharper within-cohort profiling, and that is a natural target for future work (\S6.5).

\section{Limitations of the Study}

The framework rests on a small number of design choices and data-collection realities whose implications run through every result in the thesis. Three sets of limitations stand out.

The first set concerns the data and sample. The corpus is drawn exclusively from newspaper-reported completed suicides, which biases coverage towards newsworthy cases --- unusual methods, public locations, prominent individuals --- and under-represents remote rural deaths and stigmatised populations; the findings therefore describe patterns within reported cases, not the true epidemiological distribution of suicide in Bangladesh. At 1,617 records, the corpus is, to the best of our knowledge at the time of writing (mid-2026), the largest manually curated Bangladeshi newspaper-based suicide dataset publicly assembled, but still small by machine-learning standards, and the heavy missingness further reduces effective information content, limiting cluster-boundary resolution and statistical power within clusters. A larger dataset would also enable three-or-more-cluster solutions that the current sample size leaves under-resolved.

The second set concerns methodological scope. The study is fully unsupervised, so there is no ground-truth label set against which to validate cluster assignments; the temporal coherence of the best-pipeline clusters supplies indirect support, but external validation against clinical records or police reports would be more decisive. Hyperparameter exploration within each algorithm was bounded rather than exhaustive: only a subset of DBSCAN $(\varepsilon,\,\text{minSamples})$ combinations was tested, and K-Means and Agglomerative were swept over $k = 2$--$20$ and $k = 2$--$10$ respectively, but $k = 2$ consistently won and higher-$k$ structures were not deeply probed. Two further constraints affected the feature space: the Bangla free-text columns (\texttt{profession\_description}, \texttt{reason\_description}) were excluded because no project-tier Bangla NLP was available, and weather variables remained substantially incomplete for rural locations because of sparse Meteostat station coverage. Both gaps likely understate the contribution of motivation-related and environmental features to the recovered structure.

The third set concerns two confounds that carry through the entire analysis and need separate emphasis. \phantomsection\label{lim:source_cohort}The first is source-cohort confounding in the cluster structure: as documented in Finding~4, the two-cluster solution from the best pipeline aligns almost perfectly with the dataset's two-cohort construction history --- the 2020-focused Kaggle subset versus the manually collected 2017--2019 and 2021--2025 cases --- and the two cohorts differ systematically in completeness, reporting era, and the journalistic and geographic profile of the underlying news coverage. Any cluster-level contrast in demographic, environmental, or behavioural features reported in Chapter~4 is therefore partially confounded with collection-source heterogeneity rather than cleanly isolating substantive variation, and the cluster profiles should be treated as exploratory subgroup analysis whose findings need to be re-examined against more homogeneously sourced future data. \phantomsection\label{lim:temp_units}The second is a temperature-unit inconsistency across cohorts: the raw temperature variables (\texttt{feels\_like}, \texttt{temp\_min}, \texttt{temp\_max}) inherit Kelvin units from the Kaggle 2020 subset and Celsius units from the manually collected 2017--2025 cases (filled via Meteostat in this study), and the units were not harmonised before imputation and clustering because re-running the full 32-pipeline experiment was not feasible within the project timeline. Cluster-level differences in temperature-derived features therefore partly reflect the unit mismatch rather than substantive environmental contrast, and any environment-based interpretation of the cluster split should be read with strong caution; replications of this work should harmonise temperature units before imputation.

\section{Directions for Future Work}

The findings and limitations above point toward two interconnected lines of follow-up work.

The first line aims to sharpen the substantive reading of the recovered clusters. An immediate next step is to investigate which specific features, beyond cohort membership, drive the two-cluster split; feature-importance analysis --- random forests trained on the cluster labels, or SHAP attributions over the best pipeline's representation --- would help separate genuine sociological shifts across time periods from reporting-practice artefacts and would directly inform how prevention strategies should adapt to Bangladesh's evolving social and economic landscape. Parallel data work should focus on vernacular Bangla newspapers (currently under-represented) and on more consistent attribute completion across years, so that three-or-more cluster solutions become tractable and student, rural-male, and urban-domestic-conflict risk profiles can potentially separate as distinct interpretable groups rather than within-cluster variation. Closing the Bangla NLP gap is a related, high-leverage direction: applying Bangla-specific transformer models such as BanglaBERT \cite{ref47} to \texttt{reason\_description} and \texttt{profession\_description} would recover information currently dropped from the feature space and enable finer characterisation of motivation, occupation, and social context.

The second line aims to move from descriptive subgroup analysis towards actionable prevention. The cluster-labelled dataset opens a path to causal modelling --- propensity-score matching, structural equation modelling, and related techniques \cite{ref48} --- which would test how socioeconomic stressors, seasonal factors, and demographic position translate into elevated risk within each cluster. Longer term, the cluster profiles could underpin a lightweight near-real-time alerting framework that flags elevated-risk geographies or time periods to mental-health services and NGOs whenever monitorable features (geography, time of year, occupation, family circumstances) align with high-risk-cluster patterns; this applied extension would deliver on the study's beneficence objective and align with the AI-assisted suicide-prevention literature \cite{ref43,ref44,ref53}. None of the pipeline components --- API-based conditional filling, multi-method imputation comparison, balance-aware composite evaluation, and interpretable cluster profiling --- is specific to Bangladesh or to suicide data, so validating and extending the framework on comparable severely incomplete newspaper-derived datasets in other South Asian or Sub-Saharan African contexts \cite{ref23} would help build a generalisable toolkit for exploratory subgroup analysis on incomplete social data in resource-constrained settings.

\section{Concluding Remarks}

The thesis has shown that interpretable, structurally coherent groupings can be recovered from severely incomplete, newspaper-derived social data --- even when missingness exceeds 90\% for several key features --- when imputation, clustering, and evaluation are matched to the demands of the data type. The 32-pipeline comparison, anchored by a balance-aware composite score, is offered as a starting point rather than a definitive analytical resolution.

The two-cluster solution from the best pipeline --- Agglomerative Clustering with MissForest no-context imputation --- emerged without any temporal or source-based supervision. As Finding~4 and Limitation~\ref{lim:source_cohort} argued, the boundary it recovered is best read as a partial confound between substantive variation and collection-source heterogeneity. Within those bounds the result is useful in two specific ways: it acts as an internal consistency check on the imputation pipeline, which has not collapsed cohort-level structure into a single homogeneous distribution, and it surfaces a small set of feature-level contrasts that more homogeneously sourced future data can re-examine.

The motivation behind the work is not methodological novelty alone. It is the hope that better analytical tools, applied to existing --- however imperfect --- data sources, can contribute to a more evidence-based approach to suicide prevention in Bangladesh. Every cluster profile, geographic pattern, and temporal trend reported here should be treated as a candidate hypothesis for follow-up study rather than a deployment-ready target. The 10,000 or more lives lost annually to suicide in Bangladesh \cite{ref1,ref50} are the reason this kind of careful, transparent, ethically grounded work matters.
